\documentclass[11pt,addpoints]{exam}   	% use "amsart" instead of "article" for AMSLaTeX format
\usepackage{geometry}                		% See geometry.pdf to learn the layout options. There are lots.
\geometry{letterpaper}                   		% ... or a4paper or a5paper or ... 
%\geometry{landscape}                		% Activate for rotated page geometry
%\usepackage[parfill]{parskip}    		% Activate to begin paragraphs with an empty line rather than an indent
\usepackage{graphicx}				% Use pdf, png, jpg, or eps§ with pdflatex; use eps in DVI mode
								% TeX will automatically convert eps --> pdf in pdflatex		
\usepackage{amsmath}
\usepackage{cancel}
\usepackage{amssymb}
\usepackage{multicol}
\usepackage{mhchem}
\usepackage{lewis}

\title{Quarter 3 Exam}
\author{Mark Peever}
\date{March 6, 2026}							% Activate to display a given date or no date

\begin{document}
\maketitle

\pointsinrightmargin
%\marginpointname{ \points}
\printanswers

\begin{center}
\fbox{\fbox{\parbox{5.5in}{\centering
Answer all questions on a separate sheet of paper. Show all your work.
}}}
\end{center}
\vspace{0.1in}
\makebox[\textwidth]{Name:\enspace\hrulefill}
\vspace{0.2in}

\begin{questions}

\question[5] Am\'elie and Avery are fascinated by \ce{HCl}. Am\'elie is particularly interested in the electron configurations for the atoms in it. 
She's certain the electron configuration for \ce{H} is $1s^{1}$, and she's pretty sure she can figure out the electron configuration for \ce{Cl}.
What is the electron configuration for \ce{Cl}?

\begin{solution}
$1s^{2} 2s^{2} 2p^{6} 3s^{2} 3p^{5}$
\end{solution}

\vspace{0.2in}


\question[5] Avery knows \ce{HCl} is a \emph{polar covalent} compound, and she wants to write down the Lewis Structure for it.
What is the correct Lewis Structure for \ce{HCl}?

\begin{solution}
\lewis{H}{}{}{}{}{}{}{}{}\hspace{-0.5em}--\hspace{-0.75em}\lewis{Cl}{}{}{.}{.}{.}{.}{.}{.}
\end{solution}

\vspace{0.2in}

\question[5] No\"emi shares the obsession with \ce{HCl}. She knows it's an acid, which means it produces \ce{H3O^{+}} ions in water.
What is the correct chemical equation to show \ce{HCl} producing \ce{H3O^{+}} ions in water?

\begin{solution}
\ce{HCl (aq) + H2O (l) -> H3O^{+} (aq) + Cl^{-} (aq)}
\end{solution}

\vspace{0.2in}


\question[5] No\"emi helps Thomas mix up ``stock solution" of \ce{HCl} in water, producing $1.000 L$ of $3.500 M$ \ce{HCl} solution.
How many \emph{moles} of \ce{HCl} does Thomas have?

\begin{solution}
\begin{equation} 
\begin{split}
    [HCl] &= 3.500 M \\
    V_{HCl} &= 1.000 L\\
    n &= ?\\
    [HCl] = \frac{n_{HCl}}{V_{HCl}} \\
    n_{HCl} &= [HCl] \cdot V_{HCl} \\
                 &= (3.500 \frac{mol}{L})(1.000 L) \\
                 &= 3.500 mol \\
 \end{split}
 \end{equation}
\end{solution}

\vspace{0.2in}


\question[5] Ella is curious what the \emph{mass} of Thomas's \ce{HCl} is. Since she already knows how many moles of \ce{HCl} Thomas has, it's a pretty easy calculation.
Ella knows the first step is to find the \emph{molar mass} of \ce{HCl}.
What is the molar mass of \ce{HCl}?

\begin{solution}
\begin{equation} 
\begin{split}
    M_{HCl} &= M_{H} + M_{Cl} \\
                  &= (1.00794 \frac{g}{mol}) + (35.453 \frac{g}{mol}) \\
                  &= 35.46094 \frac{g}{mol} \\
                  &= 35.461 \frac{g}{mol} \\
 \end{split}
 \end{equation}
\end{solution}

\vspace{0.2in}

\question[5] Jack, like Ella, wants to know the \emph{mass} of Thomas's \ce{HCl}. He uses the molar mass Ella calculated to calculate the mass of Thomas's \ce{HCl}.
How many \emph{grams} of \ce{HCl} does Thomas have?

\begin{solution}
\begin{equation} 
\begin{split}
    m &= n_{HCl} \cdot M_{HCl} \\
        &= (3.500 mol) (35.461 \frac{g}{mol}) \\
        &= 124.11329 mol \\
        &= 124.1 mol
 \end{split}
 \end{equation}
\end{solution}

\vspace{0.2in}

%\clearpage
\question[5] Vianne and Ezra have an Erlenmeyer containing \ce{NaOH}. Their flask contains $1.000 L$ of \ce{NaOH} solution, but they can't read Mr. Peever's handwriting, and thus don't know what the concentration is. They know \ce{NaOH} is an ionic compound, so they can figure out what the chemical reaction for its ionization looks like.
 Write down the chemical equation to show how \ce{NaOH} ionizes in water.

\begin{solution}
%\lewis{\ce{Na^{+}}}{}{}{}{}{}{}{}{} \lewis{\ce{Cl^{-}}}{.}{.}{.}{.}{.}{.}{.}{.}
\ce{NaOH (aq) + H2O (l) -> Na^{+} (aq) + OH^{-} (aq) + H2O (l)}
\end{solution}

\vspace{0.2in}

\question[5] Vianne and Ezra know they can use titration to figure out \ce{[NaOH]} in their flask. They can mix a small amount of Thomas's \ce{HCl} solution to a small amount of their \ce{NaOH} solution until the solution is neutral. They know the reaction will produce \ce{NaCl}, but they'll need a balanced chemical equation to work with.
Write down the balanced chemical reaction that will occur.

\begin{solution}
\ce{HCl (aq) + NaOH (aq) -> NaCl (aq) + H2O (l)}
\end{solution}

\vspace{0.2in}

\question[5] Gideon points out that it'll be easier to titrate \ce{NaOH} with \ce{HCl} if they use a weaker solution. He suggests they start with $10.00 mL$ of Thomas's solution and dilute it to $0.1000 M$.
What is the final volume of $0.1000 M$ solution they'll get?

\begin{solution}
\begin{equation} 
\begin{split}
   V_1 [HCl]_1 &= V_2 [HCl]_2 \\
   V_2 &= \frac{V_1 [HCl]_1} {[HCl]_2} \\
          &= \frac{(10.00 mL)(3.500 \frac{mol}{L} }{0.1000 \frac{mol}{L}} \\
          &= 350.0 mL
 \end{split}
 \end{equation}
\end{solution}

\vspace{0.2in}

\question[5] Caroline adds a small amount of boiled cabbage water (an indicator) to a $20.00 mL$ sample of the \ce{NaOH}. 
(She manages not to retch, thus earning everyone's admiration.) 
She slowly adds some of the diluted ($0.1000 M$) \ce{HCl} solution to the $20.00 mL$ \ce{NaOH} sample. 
She finds it takes $100.00 mL$ of \ce{HCl} to turn the indicator pink, indicating a neutral solution.
How many \emph{moles} of \ce{HCl} did Caroline add?

\begin{solution}
\begin{equation} 
\begin{split}
   [HCl] &= \frac{n}{V} \\
        n  &= ([HCl])(V) \\
            &= (0.1000 \frac{mol}{L}) (100.00 mL) (\frac{1 L}{1000 mL}) \\
            &= 0.01000 mol 
 \end{split}
 \end{equation}
\end{solution}

\vspace{0.2in}

\question[5] Nate has all the information he needs to calculate \ce{[NaOH]} for the $20.00 mL$ sample and knocks it out in only three lines. 
What is the concentration of the \ce{NaOH} in the $20.00 mL$ sample?
(It might take you a couple more lines than it takes Nate.)

\begin{solution}
\begin{equation} 
\begin{split}
   [NaOH] &= \frac{n}{V} \\
                &= \frac{0.01000 mol}{(20.00 mL)(\frac{1 L}{1000 mL})} \\
                &= 0.5000 \frac{mol}{L}
 \end{split}
 \end{equation}
\end{solution}

\vspace{0.2in}
 
\bonusquestion \emph{Bonus Question} 
There are three people in Scripture called ``son of man." Who are they?
\begin{solution}
\begin{itemize}
\item Ezekiel
\item Daniel
\item Christ
\end{itemize}
\end{solution}

\end{questions}

\begin{center}
This exam has \numquestions\ questions for a total of \numpoints\ points and 1 bonus point.
\end{center}

\pagebreak
\section*{Fun Facts}
\begin{multicols}{2}
\begin{itemize}
\item 
$ [\hspace{0.2 in}] = \frac{n}{V}$ 
\vspace{0.2in}

\item 
$ [\hspace{0.2 in}]_1 \cdot V_1=  [\hspace{0.2 in}]_2 \cdot V_2$
\vspace{0.2in}

\item 
$ 1 Molar = \frac{1 mole}{1 L}$
\vspace{0.2in}

\item 
$ 1 molal = \frac{1 mole}{1 kg}$
\vspace{0.2in}

\item 
$P V = n R T$
\vspace{0.2in}

\item 
$P_1 V_1 = P_2 V_2$
\vspace{0.2in}

\item 
$\frac{V_1}{T_1} = \frac{V_2}{T_2}$
\vspace{0.2in}

\item
$q = m c \Delta T$
\vspace{0.2in}

\item
$\Delta T = T_{final} - T_{initial}$
\vspace{0.2in}

\item
$\rho = \frac{m}{V} $
\vspace{0.2in}

\item
$T_K = T_C + 273.15 K$
\vspace{0.2in}

\item
$T_F = \frac{9}{5}T_C + 32^{\circ}F$
\vspace{0.2in}

\item
$T_C = \frac{5}{9}(T_F - 32^{\circ}F)$
\vspace{0.2in}

%\item
%$c_{water} = 1.000 \frac{calorie}{g \cdot ^{\circ}C}$
%\vspace{0.2in}
%
%\item
%$c_{water} = 4.184 \frac{J}{g \cdot ^{\circ}C}$
%\vspace{0.2in}
%
%\item
%$c_{copper} = 0.3851 \frac{J}{g \cdot ^{\circ}C}$
%\vspace{0.2in}

\item
$\rho_{water} = 1.000 \frac{g}{cm^3}$
\vspace{0.2in}

\item
$\rho_{copper} = 8.96 \frac{g}{cm^3}$
\vspace{0.2in}

\item
$ 1.0000 cm^3 = 1.0000 mL $
\vspace{0.2in}

\item
$ 1.000 cal = 4.184 J $
\vspace{0.2in}

\item
$m = n \times M$
\vspace{0.2in}

\item
$m = n \times m_{molar}$
\vspace{0.2in}

\item
$1 amu = 1 \frac{g}{mol}$ 
\vspace{0.2in}

\item
$N_A = 6.02214076 \times 10^{23} $ 
\vspace{0.2in}

\item 
$R = 0.0821 \frac{L \cdot atm}{mole \cdot K}$
\vspace{0.2in}

\end{itemize}
\vspace{0.2in}

\end{multicols}




\pagebreak

\begin{table}[h]
\centering
\begin{tabular}[p]{| l | l |}
\hline
\textbf{Ion Name} & \textbf{Formula} \\
\hline
ammonium  & \ce{NH4+} \\
hydronium  & \ce{H3O+} \\
hydroxide    & \ce{OH-} \\
chlorate       & \ce{ClO3-} \\
chlorite        & \ce{ClO2-} \\
nitrate          & \ce{NO3-} \\
nitrite           & \ce{NO2-} \\
acetate        & \ce{C2H3O2-} \\
cyanide       & \ce{CN-} \\
carbonate   & \ce{CO3^{2-}} \\
chromate    & \ce{CrO4^{2-}} \\
dichromate & \ce{Cr2O7^{2-}}\\
sulfate        & \ce{SO4^{2-}} \\
sulfite         & \ce{SO3^{2-}} \\
phosphate  & \ce{PO4^{3-}} 
\end{tabular}
\caption{Polyatomic Ions}
\label{table:polyatomic-ions}
\end{table}

\begin{table}
\centering
\begin{tabular}[p]{| l | l |  l |}
\hline
\textbf{Substance} & \textbf{Specific Heat ($\frac{J}{g \cdot ^{\circ}C}$)} & \textbf{Specific Heat ($\frac{calorie}{g \cdot ^{\circ}C}$)}  \\
\hline
water        & $4.184 $     & $1.000 $ \\
copper     & $0.3851$     & $0.0920 $  \\
iron          & $0.4521 $    & $0.1081 $ \\
glass        & $0.8372 $    & $0.2001 $\\
aluminum & $0.9000 $   & $0.2151 $ \\
\end{tabular}
\caption{Specific Heats}
\label{table:covalent-prefix}
\end{table}


\begin{table}
\centering
\begin{tabular}[p]{| l | l |}
\hline
\textbf{Prefix} & \textbf{Count} \\
\hline
mono & one \\
di       & two \\
tri       & three \\
tetra       & four \\
penta       & five \\
hecta       & six \\
hepta       & seven \\
octa       & eight \\
nona       & nine \\
deca       & ten \\
\end{tabular}
\caption{Prefixes for Covalent Compound Names}
\label{table:covalent-prefix}
\end{table}






\end{document}  